\documentclass[11pt]{article}
\usepackage{amsmath,amsfonts,amssymb,amsthm}
\usepackage{graphicx}
\usepackage{algorithm}
\usepackage{algorithmic}
\usepackage{float}
\usepackage{hyperref}
\usepackage{xcolor}

\newtheorem{theorem}{Theorem}
\newtheorem{lemma}{Lemma}
\newtheorem{assumption}{Assumption}

\title{Higher-order multi-scale randomized neural network (HOMS-RNN) method for high-accuracy and efficient heat conduction simulation of composite materials}
\author{Jiale Linghu, Hao Dong, Fei Wang, Yufeng Nie, Junzhi Cui}
\date{}

\begin{document}
\maketitle

\section{Introduction}

Heat conduction simulation in composite materials poses significant challenges due to highly heterogeneous microstructures with discontinuous and high-contrast material parameters. These materials exhibit complex multi-scale behaviors that are difficult to simulate using traditional numerical methods, which often require prohibitive computational resources for full fine-scale computation.

While neural network-based methods show promise for solving partial differential equations (PDEs), they face two major challenges when applied to multi-scale problems:
\begin{itemize}
    \item \textbf{Frequency Principle:} Deep neural networks fit objective functions from low to high frequencies and exhibit slower convergence for higher frequencies, making it difficult to capture the high-frequency oscillatory behavior of multi-scale problems.
    \item \textbf{Prohibitive Computation:} Accurately capturing the locally oscillatory response of heterogeneous materials requires a tremendous number of sampling points, leading to significant computational challenges and potential convergence issues.
\end{itemize}

This paper introduces the Higher-Order Multi-Scale Randomized Neural Network (HOMS-RNN) method, which combines the advantages of higher-order multi-scale approaches with the efficiency of randomized neural networks to address these challenges. The method is specifically designed for transient thermal problems of composite materials with highly discontinuous and high-contrast material parameters.

\section{Methodology}

\subsection{Multi-scale Heat Conduction Problem}

The transient multi-scale heat conduction problem in composite materials is described by the following parabolic equation with periodically oscillatory and highly discontinuous coefficients:

\begin{equation}
\begin{cases}
\rho^{\varepsilon}(x)c^{\varepsilon}(x)\frac{\partial T^{\varepsilon}(x,t)}{\partial t}-\frac{\partial}{\partial x_i}\left(k^{\varepsilon}_{ij}(x)\frac{\partial T^{\varepsilon}(x,t)}{\partial x_j}\right) = f(x,t), & \text{in } \Omega \times (0,T^*) \\
T^{\varepsilon}(x,t) = \hat{T}(x,t), & \text{on } \partial\Omega \times (0,T^*) \\
T^{\varepsilon}(x,0) = \tilde{T}(x), & \text{in } \Omega
\end{cases}
\end{equation}

where $T^{\varepsilon}(x,t)$ is the temperature, $\rho^{\varepsilon}(x)$ is the mass density, $c^{\varepsilon}(x)$ is the specific heat, $k^{\varepsilon}_{ij}(x)$ represents the thermal conductivity coefficient, and $f(x,t)$ denotes the internal heat source. The superscript $\varepsilon$ indicates the characteristic size of the microscopic unit cell, which causes the material parameters to be highly oscillatory and discontinuous.

\subsection{Higher-Order Multi-Scale Approach}

The HOMS-RNN method decouples the original multi-scale problem into simpler macroscopic and microscopic components through higher-order multi-scale analysis:

\begin{enumerate}
    \item \textbf{Asymptotic Expansion}: The temperature solution is expanded as:
    \begin{equation}
    T^{\varepsilon}(x,t) = T^{(0)}(x,y,t) + \varepsilon T^{(1)}(x,y,t) + \varepsilon^2 T^{(2)}(x,y,t) + O(\varepsilon^3)
    \end{equation}
    where $y = x/\varepsilon$ is the microscopic variable.
    
    \item \textbf{Lower-Order Cell Problem}: Solving for $T^{(1)}$ leads to:
    \begin{equation}
    T^{(1)}(x,y,t) = M^{\alpha_1}(y)\frac{\partial T^{(0)}(x,t)}{\partial x_{\alpha_1}}
    \end{equation}
    where $M^{\alpha_1}(y)$ is the first-order (lower-order) auxiliary unit cell function defined by:
    \begin{equation}
    \begin{cases}
    \frac{\partial}{\partial y_i}\left(k_{ij}(y)\frac{\partial M^{\alpha_1}(y)}{\partial y_j}\right) = -\frac{\partial k_{i\alpha_1}(y)}{\partial y_i}, & \text{in } Y \\
    M^{\alpha_1}(y) = 0, & \text{on } \partial Y
    \end{cases}
    \end{equation}
    
    \item \textbf{Macroscopic Homogenized Problem}: The zeroth-order term $T^{(0)}$ is governed by:
    \begin{equation}
    \begin{cases}
    S^*\frac{\partial T^{(0)}(x,t)}{\partial t} - \frac{\partial}{\partial x_i}\left(k^*_{ij}\frac{\partial T^{(0)}(x,t)}{\partial x_j}\right) = f(x,t), & \text{in } \Omega \times (0,T^*) \\
    T^{(0)}(x,t) = \hat{T}(x,t), & \text{on } \partial\Omega \times (0,T^*) \\
    T^{(0)}(x,0) = \tilde{T}(x), & \text{in } \Omega
    \end{cases}
    \end{equation}
    where the homogenized material parameters $S^*$ and $k^*_{ij}$ are calculated as:
    \begin{equation}
    \begin{aligned}
    S^* &= \frac{1}{|Y|}\int_Y \rho(y)c(y)dY \\
    k^*_{ij} &= \frac{1}{|Y|}\int_Y \left(k_{ij}(y) + k_{ik}(y)\frac{\partial M^j(y)}{\partial y_k}\right)dY
    \end{aligned}
    \end{equation}
    
    \item \textbf{Higher-Order Cell Problems}: The second-order term $T^{(2)}$ is defined as:
    \begin{equation}
    T^{(2)}(x,y,t) = M^{\alpha_1\alpha_2}(y)\frac{\partial^2 T^{(0)}(x,t)}{\partial x_{\alpha_1}\partial x_{\alpha_2}} + R(y)\frac{\partial T^{(0)}(x,t)}{\partial t}
    \end{equation}
    where $M^{\alpha_1\alpha_2}(y)$ and $R(y)$ are second-order (higher-order) auxiliary unit cell functions defined by:
    \begin{equation}
    \begin{cases}
    \frac{\partial}{\partial y_i}\left(k_{ij}(y)\frac{\partial M^{\alpha_1\alpha_2}(y)}{\partial y_j}\right) = k^*_{\alpha_1\alpha_2} - k_{\alpha_1\alpha_2}(y) - k_{\alpha_1 j}(y)\frac{\partial M^{\alpha_2}(y)}{\partial y_j} - \frac{\partial}{\partial y_i}\left(k_{i\alpha_2}(y)M^{\alpha_1}(y)\right), & \text{in } Y \\
    M^{\alpha_1\alpha_2}(y) = 0, & \text{on } \partial Y
    \end{cases}
    \end{equation}
    \begin{equation}
    \begin{cases}
    \frac{\partial}{\partial y_i}\left(k_{ij}(y)\frac{\partial R(y)}{\partial y_j}\right) = \rho(y)c(y) - S^*, & \text{in } Y \\
    R(y) = 0, & \text{on } \partial Y
    \end{cases}
    \end{equation}
    
    \item \textbf{Higher-Order Multi-Scale Solution}: The final solution is assembled as:
    \begin{equation}
    T^{(2\varepsilon)}(x,t) = T^{(0)}(x,t) + \varepsilon M^{\alpha_1}(y)\frac{\partial T^{(0)}(x,t)}{\partial x_{\alpha_1}} + \varepsilon^2\left[M^{\alpha_1\alpha_2}(y)\frac{\partial^2 T^{(0)}(x,t)}{\partial x_{\alpha_1}\partial x_{\alpha_2}} + R(y)\frac{\partial T^{(0)}(x,t)}{\partial t}\right]
    \end{equation}
\end{enumerate}

This approach effectively decomposes the multi-scale solution into non-oscillatory macroscopic behavior, low-frequency oscillatory information, and high-frequency oscillatory information.

\subsection{Randomized Neural Networks}

The HOMS-RNN method employs randomized neural networks (RNNs) to efficiently solve the separated problems:

\begin{enumerate}
    \item \textbf{RNN Structure}: A fully connected neural network $\mathcal{N}(x;\theta)$ with depth $D$ is defined as:
    \begin{equation}
    \begin{aligned}
    \varphi_0(x) &= x \\
    \varphi_l(x) &= \sigma(W_l\varphi_{l-1} + b_l), \quad l = 1,\ldots,D-1 \\
    \mathcal{N}(x;\theta) &= \varphi_D(x) = W_D\varphi_{D-1} + b_D
    \end{aligned}
    \end{equation}
    where $\sigma(\cdot)$ is the activation function (typically tanh).
    
    \item \textbf{Randomization}: In an RNN, all biases and weights $[(W_1,b_1),\ldots,(W_{D-1},b_{D-1})]$ are randomly initialized from $[-R_m,R_m]$ and fixed (non-trainable). Only the weights of the last layer $W_D$ are computed by solving a least-squares problem.
    
    \item \textbf{Function Representation}: Any function $u_\sigma \in \mathcal{N}^D_\sigma$ (one neuron in the output layer) can be expressed as:
    \begin{equation}
    u_\sigma(x) = \sum_{i=1}^{n_{D-1}} \beta_i \varphi_i(x)
    \end{equation}
    where $\varphi_i$ is the $i$-th component of the vector $\varphi_{D-1}$ and is fixed once the hidden-layer coefficients are randomly assigned.
\end{enumerate}

\subsection{HOMS-RNN Implementation}

The HOMS-RNN method uses specialized RNNs to solve each part of the multi-scale problem:

\begin{enumerate}
    \item \textbf{Microscopic Cell RNNs}: Two RNNs $u^1_\sigma(y)$ and $\tilde{u}^1_\sigma(y)$ are used to approximate the microscopic cell function $M^{\alpha_1}(y)$ in the matrix material $Y_1$ and inclusion material $Y_2$:
    \begin{equation}
    u^1_\sigma(y) = \sum_{m=1}^{M_1} \beta^1_m \varphi^1_m(y), \quad \tilde{u}^1_\sigma(y) = \sum_{m=1}^{M_2} \tilde{\beta}^1_m \tilde{\varphi}^1_m(y)
    \end{equation}
    
    \item \textbf{Hard-Constraint Boundary Conditions}: To improve accuracy, hard-constraint boundary conditions are applied:
    \begin{equation}
    \hat{u}^1_\sigma(y) = h(y) + \ell(y)u^1_\sigma(y)
    \end{equation}
    where $h(y) = 0$ for homogeneous Dirichlet boundary conditions and $\ell(y)$ is a function that equals zero on the boundary and is positive inside the domain.
    
    \item \textbf{Space-Time RNN}: A space-time RNN $u^0_\sigma(x,t)$ approximates the macroscopic homogenized solution $T^{(0)}(x,t)$:
    \begin{equation}
    u^0_\sigma(x,t) = \sum_{m=1}^{M} \beta^0_m \varphi^0_m(x,t)
    \end{equation}
    
    \item \textbf{Higher-Order RNNs}: Similar RNNs are used for the higher-order cell functions $M^{\alpha_1\alpha_2}(y)$ and $R(y)$.
    
    \item \textbf{Automatic Differentiation Assembly}: Once all components are computed, automatic differentiation is used to assemble the final HOMS-RNN solution:
    \begin{equation}
    \hat{T}^{(2\varepsilon)}(x,t) = \hat{T}^{(0)}(x,t) + \varepsilon \hat{M}^{\alpha_1}(y)\frac{\partial \hat{T}^{(0)}(x,t)}{\partial x_{\alpha_1}} + \varepsilon^2\left[\hat{M}^{\alpha_1\alpha_2}(y)\frac{\partial^2 \hat{T}^{(0)}(x,t)}{\partial x_{\alpha_1}\partial x_{\alpha_2}} + \hat{R}(y)\frac{\partial \hat{T}^{(0)}(x,t)}{\partial t}\right]
    \end{equation}
\end{enumerate}

The algorithm can be summarized in the following steps:
\begin{enumerate}
    \item Decouple original multi-scale PDEs into macroscopic homogenized PDEs and microscopic cell PDEs using the higher-order multi-scale model.
    \item Solve lower-order microscopic cell functions using RNNs with hard-constraint boundary conditions.
    \item Calculate macroscopic homogenized material parameters using the obtained microscopic cell functions.
    \item Solve the macroscopic homogenized problem using a space-time RNN.
    \item Solve higher-order microscopic cell functions using RNNs.
    \item Assemble the final solution using automatic differentiation.
\end{enumerate}

\section{Theoretical Findings}

The paper provides a rigorous error analysis of the HOMS-RNN method. The key theoretical results include:

\begin{lemma}
Based on universal approximation theorem, for any $u \in C(I^d)$ with $I^d = [a,b]^d$ $(d = 1,2,3,\ldots)$ and $\forall \varepsilon > 0$, there exists a function $u_\sigma \in \mathcal{N}^D_\sigma$, as long as the hidden layer size is large enough, which is defined on any compact $K$ of $C(I^d)$. Then it follows that
\begin{equation}
\sup_{x \in K} |u_\sigma(x) - u(x)| < \varepsilon
\end{equation}
\end{lemma}

\begin{assumption}
For any $u \in C^3(I^d)$ and $\forall \varepsilon > 0$, there exists a function $u_\sigma \in \mathcal{N}^D_\sigma$, as long as the hidden layer size is large enough, which is defined on any compact $K$ of $C^3(I^d)$. Then we have
\begin{equation}
\max_{|\alpha| \leq 3} \sup_{x \in K} |\partial^\alpha_x u_\sigma(x) - \partial^\alpha_x u(x)| < \varepsilon
\end{equation}
\end{assumption}

\begin{lemma}
According to results in the literature, the error estimation for higher-order multi-scale solution $T^{(2\varepsilon)}(x,t)$ and exact solution $T^\varepsilon(x,t)$ of multi-scale parabolic problem is obtained as follows:
\begin{equation}
\sup_{0 \leq t \leq T^*} \int_\Omega \left(T^\varepsilon(x,t) - T^{(2\varepsilon)}(x,t)\right)^2 dx + \int_{0}^{T^*} \|T^\varepsilon(x,t) - T^{(2\varepsilon)}(x,t)\|^2_{H^1(\Omega)} dt \leq C(T^*)\varepsilon^2
\end{equation}
where $C(T^*)$ is a positive constant independent of $\varepsilon$, but dependent on total time $T^*$ and $\Omega$.
\end{lemma}

\begin{theorem}
Based on the above lemmas and assumption, the following global error estimation is derived for the HOMS-RNN method:
\begin{equation}
\sup_{0 \leq t \leq T^*} \int_\Omega \left(T^\varepsilon(x,t) - \hat{T}^{(2\varepsilon)}(x,t)\right)^2 dx + \int_{0}^{T^*} \|T^\varepsilon(x,t) - \hat{T}^{(2\varepsilon)}(x,t)\|^2_{H^1(\Omega)} dt \leq C(T^*)\varepsilon^2 + C(T^*)\varepsilon^2
\end{equation}
where $C$ is a positive constant independent of $\varepsilon$ and $\varepsilon$, but dependent on total time $T^*$ and $\Omega$.
\end{theorem}

This theorem establishes the convergence of the HOMS-RNN method, showing that the error consists of two parts: the modeling error from the higher-order multi-scale approximation and the approximation error from the randomized neural networks.

\section{Numerical Examples and Results}

The paper presents several numerical experiments to demonstrate the effectiveness of the HOMS-RNN method:

\subsection{Drawbacks of Direct RNN Simulation}

The authors first demonstrate why direct RNN simulation is ineffective for multi-scale problems:

\begin{itemize}
    \item As the scale parameter $\varepsilon$ decreases (e.g., $1, 1/2, 1/3, 1/4, 1/5$), direct RNN simulations fail to accurately capture the oscillatory behavior at the micro-scale.
    \item Computational time increases dramatically as $\varepsilon$ decreases ($2.33$s for $\varepsilon=1$ to $1830.26$s for $\varepsilon=1/5$).
    \item Relative errors in direct RNN simulations are unacceptably high, especially for smaller values of $\varepsilon$.
\end{itemize}

\subsection{Composite Materials with Different Thermal Conductivities}

The HOMS-RNN method was tested on composite materials with different thermal conductivity contrasts (10:1, 50:1, and 100:1):

\begin{itemize}
    \item \textbf{Network Configuration}: RNNs with 320 neurons in the hidden layer for microscopic cell functions and 640 neurons for macroscopic homogenized solutions; 9,000 residual points within the microscopic cell and 3,000 points along the interface.
    
    \item \textbf{Hard vs. Soft Constraints}: Hard-constraint boundary conditions showed better performance than soft constraints, with lower relative errors for all cell functions and the macroscopic solution.
    
    \item \textbf{Results for 10:1 Contrast}: The HOMS-RNN method effectively captured the microscopic oscillatory phenomena and showed stable performance throughout the temporal evolution, with errors comparable to FEM.
    
    \item \textbf{High-Contrast Materials}: For 50:1 and 100:1 contrast ratios, the lower-order multi-scale method failed to effectively capture microscopic oscillations, while the HOMS-RNN method maintained high accuracy and stability.
\end{itemize}

\subsection{Robustness Analysis}

The method was tested with progressively smaller characteristic parameter $\varepsilon$ values (1/5, 1/10, 1/20, 1/30):

\begin{itemize}
    \item The HOMS-RNN method accurately captured the oscillatory phenomena even as the number of microscopic unit cells increased dramatically.
    
    \item Computational efficiency remained stable as $\varepsilon$ decreased, while FEM computation time increased exponentially (from 144.81s for $\varepsilon=1/5$ to 7016.75s for $\varepsilon=1/30$).
    
    \item The HOMS-RNN method took only 96.86s for $\varepsilon=1/5$ and 106.46s for $\varepsilon=1/30$, showing remarkable efficiency for large-scale simulations.
\end{itemize}

\subsection{3D Composite Materials}

The method was applied to 3D SiC/C composite materials:

\begin{itemize}
    \item \textbf{Network Configuration}: RNNs with 1,280 neurons in the hidden layer; 40,000 residual points within the microscopic unit cell and 20,000 points along the interface.
    
    \item \textbf{Results}: The HOMS-RNN method accurately captured the high-frequency oscillatory response in the 3D composite, with excellent agreement with reference FEM solutions.
    
    \item \textbf{Computational Efficiency}: The 3D simulation using HOMS-RNN required only 1.5 hours, compared to 7.1 hours for the higher-order multi-scale method and 11.8 hours for reference FEM.
\end{itemize}

\subsection{Porous Materials}

The method was also tested on two types of porous materials:

\begin{itemize}
    \item For both simple and complex porous geometries, the HOMS-RNN method achieved competitive accuracy with traditional numerical methods while maintaining superior computational efficiency.
    
    \item As $\varepsilon$ decreased, computational errors of the HOMS-RNN method converged, demonstrating the robustness of the approach.
    
    \item For complex porous material, the HOMS-RNN method took 225.81s for $\varepsilon=1/5$ and 291.85s for $\varepsilon=1/30$, while FEM required 193.97s and 17,657.60s respectively.
\end{itemize}

\section{Implications and Advantages}

The HOMS-RNN method offers several significant advantages:

\begin{enumerate}
    \item \textbf{Addressing the Frequency Principle}: By decomposing the multi-scale solution into different frequency components and using higher-order multi-scale constraints, the method effectively captures both low and high-frequency oscillatory behavior.
    
    \item \textbf{Computational Efficiency}: The method dramatically reduces computational time compared to traditional approaches, especially for problems with small characteristic scales. The efficiency remains stable as the scale parameter decreases.
    
    \item \textbf{Mesh-Free Advantage}: The method eliminates the need for mesh generation, which is particularly valuable for problems with complex geometries.
    
    \item \textbf{High Accuracy}: The method achieves accuracy comparable to fine-scale FEM simulations, particularly for high-contrast materials where classical methods struggle.
    
    \item \textbf{Space-Time Approach}: The space-time formulation for the macroscopic problem allows efficient solution of time-dependent problems without time-stepping schemes.
    
    \item \textbf{Robustness}: The method maintains its accuracy and efficiency across different material properties, geometries, and scale parameters.
\end{enumerate}

These advantages make the HOMS-RNN method particularly suitable for large-scale engineering applications involving composite materials with complex microstructures and high-contrast properties. The method's ability to efficiently handle 3D simulations and complex geometries further enhances its practical value.

Future directions for this approach include extending it to multi-physics coupling problems, leveraging parallel computing capabilities of RNNs for further efficiency improvements, and addressing nonlinear and random multi-scale problems.

\end{document} 